\documentclass[sigconf, nonacm]{acmart}

%%% do not modify the following VLDB block %%
%%% VLDB block start %%%
\usepackage{pvldb}
%%% VLDB block end %%%

\renewcommand\vldbdoi{XX.XX/XXX.XX}
\renewcommand\vldbpages{XXX--XXX}
\renewcommand\vldbavailabilityurl{}

\newcommand{\E}{\mathcal{E}}
\newcommand{\Q}{\mathcal{Q}}
\newcommand{\Pmatch}{\mathcal{P}}
\newcommand{\Rset}{\mathbb{R}}

\begin{document}
\title{Budgeted Extended Statistics Selection under Maximal-Matching Estimator Uncertainty}

\author{Anonymous Author(s)}
\affiliation{%
  \institution{Anonymous Institution}
  \city{Anonymous City}
  \country{Anonymous Country}
}

\begin{abstract}
Placeholder abstract for the problem-formulation manuscript.
\end{abstract}

\maketitle

%%% do not modify the following VLDB block %%
%%% VLDB block start %%%
\vldbtopmatter
%%% VLDB block end %%%

\section{Problem Formulation}

\subsection{Workload and Budgeted Resources}

Let \(\mathcal I\) be a finite collection of cardinality-estimation tasks.
Task \(i\in\mathcal I\) identifies a well-defined cardinality target, which
may be a selection, join result, logical intermediate, or other explicitly
specified relational expression.  It is not necessarily an entire SQL
query or a physical plan node.  Each task has weight \(w_i>0\), ground-truth
cardinality \(N_i\geq0\), and target \(T_i=\max\{1,N_i\}\).

Let \(\E\) be a finite universe of candidate estimation resources.  A
resource \(e\in\E\) has positive creation cost \(c_e\).  The advisor chooses
one global, unordered creation set \(M\subseteq\E\) satisfying
\[
  \sum_{e\in M}c_e\leq B,\qquad B\geq0.
\]
The budget is an upper bound; no monotonicity of estimation accuracy is
assumed.  The generic formulation places no restriction on a resource's
arity, representation, or implementation.

\subsection{Admissible Schemes and Estimation Models}

For each task \(i\), let \(\Omega_i\) be a finite space of complete
estimation schemes.  A scheme \(\sigma\in\Omega_i\) contains all information
needed to identify how its estimate is computed and may encode resources,
order, coverage, or other estimator-specific choices.  Let
\(\operatorname{used}(\sigma)\subseteq\E\) denote the resources required by
\(\sigma\).  The abstract admissibility mapping is
\[
  \Gamma_i:2^\E\longrightarrow
    2^{\Omega_i}\setminus\{\varnothing\},
\]
with
\[
  \sigma\in\Gamma_i(M)\ \Longrightarrow\
  \operatorname{used}(\sigma)\subseteq M.
\]
We assume \(\Gamma_i(M)\) is nonempty for every created set, including
\(M=\varnothing\), but impose no monotonicity under set inclusion:
\(M_1\subseteq M_2\) need not imply
\(\Gamma_i(M_1)\subseteq\Gamma_i(M_2)\).  Created
resources and resources used by an admissible scheme are therefore
distinct concepts; schemes need not be sets, matchings, sequences, or
physical execution plans.

Separately, let
\[
  \operatorname{Est}_i:\Omega_i\longrightarrow\mathbb{R}_{>0}
\]
be a black-box cardinality estimator.  No algebraic property is required:
the generic model assumes neither separability, independence, additive
contributions, nor a particular output conversion.  The admissibility
mapping specifies which schemes may be used, whereas \(\operatorname{Est}_i\)
specifies the estimate produced by each scheme.  For a nonnegative loss
function \(\ell_i:\Omega_i\to\mathbb{R}_{\geq0}\), we use absolute
logarithmic error as the running example:
\[
  \ell_i(\sigma)=
  \left|\log\frac{\operatorname{Est}_i(\sigma)}{T_i}\right|.
\]

\subsection{General Pessimistic Selection}

The estimator-agnostic PESS-ESS problem is
\[
\boxed{\begin{aligned}
  \operatorname*{minimize}_{M\subseteq\E}\quad
    &\sum_{i\in\mathcal I}w_i
      \max_{\sigma\in\Gamma_i(M)}\ell_i(\sigma)\\
  \text{subject to}\quad
    &\sum_{e\in M}c_e\leq B.
\end{aligned}}
\]
The created configuration determines the admissible schemes, and the
estimation model assigns loss to each scheme.  If a family of actual
policies \(\Theta\) induces \(\sigma_i(M;\theta)\) with
\(\sigma_i(M;\theta)\in\Gamma_i(M)\) for every relevant \(i,M,\theta\), then
\[
  \max_{\theta\in\Theta}\sum_{i\in\mathcal I}w_i
    \ell_i(\sigma_i(M;\theta))
  \leq
  \sum_{i\in\mathcal I}w_i
    \max_{\sigma\in\Gamma_i(M)}\ell_i(\sigma).
\]
The inequality is conditional on that inclusion and can be strict when a
single policy couples choices across tasks; it is not a universal
guarantee about any database system's realized error.

If a finite query workload \(\mathcal Q\) contributes exactly one task per
query, then
\(\mathcal I\cong\mathcal Q\) and the query-indexed formulation is recovered;
a query may instead contribute several logical cardinality tasks.  A task
specifies what cardinality is estimated, a scheme specifies how estimation
resources and decisions are used, and a physical execution plan specifies
how a query is executed.  PESS-ESS concerns the first two and does not
select plans or minimize estimated execution cost.  Its task-wise maxima are
taken independently, so they may be conservative when several tasks share a
query or implementation state.

\subsection{Pairwise Extended-Statistics Instantiation}

We now instantiate the abstract problem for the single-table conjunctive
cardinality-estimation tasks studied here, including the two-column MCV
statistics that motivate the setting.  Thus, the concrete task set
\(\mathcal I\) is restricted to these tasks; the following matching model is
not asserted for arbitrary join or intermediate-result tasks.

Each candidate covers exactly two attributes.  For each task \(i\), let
\(\E_i\subseteq\E\) be the structurally eligible candidates, based on
statistic kind, covered columns, and supported predicate operators.  They
form the undirected edge-labeled multigraph
\[
  H_i=(V_i,\E_i),\qquad V_i=\bigcup_{e\in\E_i}A_e,\qquad |A_e|=2.
\]
Each edge is a distinguishable candidate statistic; parallel edges therefore
remain distinct even when they have the same endpoints.  For a created set
\(M\), let \(H_i[M]\) contain the edges \(M\cap\E_i\).

The concrete admissibility rule is
\[
  \boxed{\Gamma_i^{\mathrm{MM}}(M)=
    \operatorname{MaximalMatchings}(H_i[M]).}
\]
A matching is a set of edges with no shared endpoint, and it is maximal if
no additional available edge can be added; maximal does not mean maximum
cardinality.  The empty graph has the unique empty matching.  This rule is
a deliberate abstraction of a restricted PostgreSQL estimator behavior
motivated by compatible clause coverage, not a universal characterization of
PostgreSQL, arbitrary predicates, or arbitrary-arity statistics.

For the concrete estimator, let \(n_i^0>0\) be the baseline pre-clamp
estimated cardinality and let \(r_{e,i}>0\) be a fixed correction factor for
each \(e\in\E_i\).  For any pairwise compatible matching \(P\) in \(H_i[M]\),
we instantiate the black-box estimator as
\[
  \boxed{\operatorname{Est}_i^{\mathrm{mult}}(P)
    =n_i^0\prod_{e\in P}r_{e,i}}.
\]
The factors are calibrated from full-precision, pre-clamp estimates or
selectivities, not rounded integer \texttt{Plan Rows}.  The assumption is
\emph{multiplicative separability}: each factor is fixed for the task and
does not depend on which other compatible resources are used.  This is not
an independence assumption about the data; a multivariate statistic may
already encode dependence within its own attribute group.

Define
\[
  \alpha_i=\log\frac{n_i^0}{T_i},\qquad
  b_{e,i}=\log r_{e,i}.
\]
The resulting raw loss is
\[
  \ell_i^{\mathrm{raw}}(P)=
  \left|\alpha_i+\sum_{e\in P}b_{e,i}\right|.
\]

\subsection{Concrete Pessimistic Selection}

Substituting \(\Gamma_i^{\mathrm{MM}}\) and the raw loss into the general
problem gives the structured raw-loss PESS-ESS instance:
\[
\boxed{\begin{aligned}
  \operatorname*{minimize}_{M\subseteq\E}\quad
    &\sum_{i\in\mathcal I}w_i
      \max_{P\in\operatorname{MaximalMatchings}(H_i[M])}
      \left|\alpha_i+\sum_{e\in P}b_{e,i}\right|\\
  \text{subject to}\quad
    &\sum_{e\in M}c_e\leq B.
\end{aligned}}
\]
The raw-loss form is the structured objective used for the present
theoretical analysis; plan loss is a separate finalized-row objective.
When finalized database row estimates are the target, the same instantiation
can instead use the plan-loss transformation
\[
  C(x)=
  \begin{cases}
    1,&0<x\leq1,\\
    \operatorname{rint}(x),&1<x\leq x_{\max},
  \end{cases}
\]
for finite positive \(x\) in the ordinary implementation regime.  Here
\(x_{\max}\) denotes the implementation's finite cap and
\(\operatorname{rint}\) follows its floating-point rounding mode; nonfinite
inputs and implementation-specific very-large-value paths are outside this
instantiation.  The plan loss is
\[
  \ell_i^{\mathrm{plan}}(P)=
  \left|\log C\!\left(\operatorname{Est}_i^{\mathrm{mult}}(P)\right)
    -\log T_i\right|.
\]
Raw and plan loss are distinct objectives and need not have the same
optimizer.  A secondary comparator may average plan loss uniformly over
\(\operatorname{MaximalMatchings}(H_i[M])\), but that distribution is not in
general the distribution induced by uniform creation permutations.

\end{document}
